\documentclass[11pt]{article}
\usepackage{amsmath}
\usepackage{amssymb}
\usepackage{amsthm}
\usepackage{enumitem}

\begin{document}
\newcommand{\seq}[2][n]{\ensuremath{\left\{ #2 _{#1} \right\}}}
\section{Theorem 16.7}
    Let $f: X \to Y$ be continuous and let $E \subset X$ be connected.
    Then $f(E) \subset Y$ is connected.
    \begin{proof}[Proof by Contradiction]
        Suppose $E \subset X$ is connected but suppose for contradiction that $f(E) \subset Y$ is not connected.
        We may write $f(E) = A \cup B$ with open, disjoine, nonempty sets $A$ and $B$ in $(f(E), d_Y)$.
        By properties of images and preimages, we have $E \subset f^{-1}(f(E))$.
        \begin{equation}
            f^{-1}(f(E)) = f^{-1}(A \cup B) = f^{-1}(A) \cup f^{-1}(B)
        \end{equation}

        Therefore, $E = (f^{-1}(A) \cap E) \cup (f^{-1}(B) \cap E)$.
        Let $C = f^{-1}(A) \cap E$ and $D = f^{-1}(B) \cap E$.

        \underline{$C$ and $D$ are disjoint.}
            If they weren't disjoint, then there would be $x$ such that $x \in f^{-1}(A)$ and $x \in f^{-1}(B)$.
            That would mean that $f(x) \in A$ and $f(x) \in B$, which would not be possible since $A$ and $B$ are disjoint.

        \underline{$C$ and $D$ are nonempty.}
            $A$ and $B$ are nonempty.
            If $C = (f^{-1}(A) \cap E)$ were empty, then $E$ would be a subset of $f^{-1}(B)$.
            This would imply $f(E) \subset B$, but we know $f(E) \subset B$, which would mean $f(E) = B$.
            That would mean $A = \varnothing$, but we already established this is false.
        
        \underline{$C$ and $D$ are open in $(E,d_X)$.}
            Since $A$ and $B$ are open, we know that for $U$ and $V$ open in $Y$, $A = U \cap f(E)$ and $B = V \cap f(E)$.
            We can therefore say the following:
            \begin{align}
                f^{-1}(A)   &=  f^{-1}(U \cap f(E)) 
                    =   f^{-1}(U) \cap f^{-1}(f(E))\\
                f^{-1}(B)   &=  f^{-1}(V \cap f(E)) 
                    =   f^{-1}(V) \cap f^{-1}(f(E))\\
                f^{-1}(A) \cap E    &=  f^{-1}(U) \cap f^{-1}(f(E)) \cap E
                    =   f^{-1}(U) \cap E\\
                f^{-1}(B) \cap E    &=  f^{-1}(V) \cap f^{-1}(f(E)) \cap E
                    =   f^{-1}(V) \cap E
            \end{align}
            Since by continuity $f^{-1}(U)$ and $f^{-1}(V)$ are open in $(X,d_X)$, $f^{-1}(A) \cap E$ is open in $(E,d_X)$ and $f^{-1}(B) \cap E$ is open in $(E,d_X)$, so $C$ and $D$ are open in $(E,d_X)$.

        Therefore, $E = C \cup D$ for open, disjoint, nonempty $C$ and $D$.
        Therefore, $E$ is separated, which is a contradiction.
        Ergo, $f(E)$ is connected.
    \end{proof}

\section{Theorem 17.7}
    Let $f: X \to Y$ be continuous between metric spaces.
    Let $X$ be compact.
    Then, $f$ is uniformly continuous.
    \begin{proof}
        Let $\varepsilon > 0$.
        For all $p \in X$, there exists $\delta_p > 0$ (depending on $p$) such that if $d_X(p,q) < \delta_p$, then $d_Y(f(p),f(q)) < \frac{\varepsilon}{2}$.
        Therefore, the open sets $\left\{ B_{\frac{\delta_p}{2}}(p) \right\}$ cover $X$.
        Since $X$ is compact, there is a finite subcover $\left\{ B_{\frac{\delta_{p_k}}{2}}(p_k) \right\}_{k=1}^n$.
        \begin{equation}\tag{$\star$} \label{eq:star}
            X \subset \left\{ B_{\frac{\delta_{p_k}}{2}}(p_k) \right\}_{k=1}^n
        \end{equation}

        Let $\delta = \min_{k \in \mathbb{N} \leq n} \left\{ \frac{\delta_{p_k}}{2} \right\} > 0$.
        Choose any $x,y \in X$ with $d_X(x,y) < \delta$.
        Since (\ref{eq:star}) is true, there exists $j$ where $1 \leq j \leq n$ and $x \in B_{\frac{\delta_{p_j}}{2}}(p_j)$.
        Therefore, we can say things about $d(x,p_j)$ and $d(y,p_j)$.
        \begin{align}
            d_X(x,p_j)  &<  \frac{\delta_{p_j}}{2}
                <   \delta_{p_j}\\
            d_X(y,p_j)  &\leq   d_X(x,y) + d_X(x,p_j)
                \leq    \frac{\delta_{p_j}}{2} + \frac{\delta_{p_j}}{2}
                =   \delta_{p_j}
        \end{align}
        
        Since $f$ is continuous, 
        \begin{equation}
            d_Y(f(x),f(y)) \leq d_Y(f(x),f(p_j)) + d_Y(f(p_j),f(y)) < \frac{\varepsilon}{2} + \frac{\varepsilon}{2} = \varepsilon.
        \end{equation}

        Therefore, $f$ is uniformly continuous.
    \end{proof}
\end{document}